\documentclass[12pt]{article}
\usepackage{amsmath, amssymb, amsthm}
\usepackage{geometry}
\usepackage{hyperref}
\usepackage{xcolor}
\usepackage{tcolorbox}
\usepackage{enumitem}
\usepackage{mdframed}

\geometry{a4paper, margin=2.5cm}

\definecolor{problemblue}{RGB}{30, 80, 160}
\definecolor{solngreen}{RGB}{20, 120, 60}
\definecolor{notegray}{RGB}{80, 80, 80}
\definecolor{boxbg}{RGB}{240, 245, 255}

\tcbuselibrary{skins, breakable}

\newtcolorbox{problembox}[1]{
  colback=boxbg, colframe=problemblue,
  fonttitle=\bfseries\color{white},
  title={#1},
  boxrule=1pt, arc=4pt
}

\newtcolorbox{answerbox}{
  colback=green!8, colframe=solngreen,
  fonttitle=\bfseries,
  title={Final Answer},
  boxrule=1pt, arc=4pt
}

\newtcolorbox{propbox}{
  colback=yellow!10, colframe=orange!70!black,
  fonttitle=\bfseries,
  title={Properties Used},
  boxrule=0.8pt, arc=4pt
}

\title{\textbf{JEE Advanced — Logarithms}}
\author{Virat Dixit}
\date{\today}

\begin{document}

\maketitle

% =====================================================================
\section{Problem 1 — IIT-JEE 2011}
% =====================================================================

\begin{problembox}{Problem 1 \hfill \normalfont\small [IIT-JEE 2011]}
Let $(x_0, y_0)$ be the solution of the following system of equations:
\begin{align}
(2x)^{\ln 2} &= (3y)^{\ln 3} \tag{i}\\
3^{\ln x} &= 2^{\ln y} \tag{ii}
\end{align}
Then $x_0$ is:
\[
\text{(A)}\ \frac{1}{6} \qquad
\text{(B)}\ \frac{1}{3} \qquad
\text{(C)}\ \frac{1}{2} \qquad
\text{(D)}\ 6
\]
\end{problembox}

\subsection*{Solution}

\textbf{Step 1: Process equation (ii).}

Take $\log$ on both sides of equation (ii):
\[
\ln x \cdot \log 3 = \ln y \cdot \log 2
\]
where all logarithms are to a common base (natural or base-10; the algebra is identical). Rearranging:
\[
\log y = \frac{\log x \cdot \log 3}{\log 2} \tag{iii}
\]

\textbf{Step 2: Process equation (i).}

Take $\log$ on both sides of equation (i):
\[
\log 2 \cdot \log(2x) = \log 3 \cdot \log(3y)
\]
Expand using the product rule $\log(ab) = \log a + \log b$:
\[
\log 2\,(\log 2 + \log x) = \log 3\,(\log 3 + \log y)
\]
\[
(\log 2)^2 + (\log 2)(\log x) = (\log 3)^2 + (\log 3)(\log y) \tag{iv}
\]

\textbf{Step 3: Substitute (iii) into (iv).}

Replace $\log y$ in equation (iv):
\[
(\log 2)^2 + (\log 2)(\log x) = (\log 3)^2 + (\log 3) \cdot \frac{\log x \cdot \log 3}{\log 2}
\]
\[
(\log 2)^2 + (\log 2)(\log x) = (\log 3)^2 + \frac{(\log 3)^2 \cdot \log x}{\log 2}
\]

Multiply throughout by $\log 2$:
\[
(\log 2)^3 + (\log 2)^2(\log x) = (\log 3)^2 \cdot \log 2 + (\log 3)^2 \cdot \log x
\]

Bring $\log x$ terms to one side:
\[
(\log 2)^2(\log x) - (\log 3)^2(\log x) = (\log 3)^2 \cdot \log 2 - (\log 2)^3
\]
\[
\log x\bigl[(\log 2)^2 - (\log 3)^2\bigr] = \log 2\bigl[(\log 3)^2 - (\log 2)^2\bigr]
\]
\[
\log x\bigl[(\log 2)^2 - (\log 3)^2\bigr] = -\log 2\bigl[(\log 2)^2 - (\log 3)^2\bigr]
\]

Dividing both sides by $[(\log 2)^2 - (\log 3)^2] \neq 0$:
\[
\log x = -\log 2 = \log\!\left(\frac{1}{2}\right)
\]
\[
\boxed{x_0 = \dfrac{1}{2}}
\]

\begin{answerbox}
\[\mathbf{x_0 = \dfrac{1}{2}}\]
Correct option: \textbf{(C)}.
\end{answerbox}

\begin{propbox}
\begin{enumerate}[leftmargin=*, label=\textbf{P\arabic*.}]
  \item \textbf{Taking logarithms of exponential equations.}  
  If $a^f = b^g$, then $f \log a = g \log b$.

  \item \textbf{Product rule of logarithms.}  
  $\log(ab) = \log a + \log b$.

  \item \textbf{Substitution to reduce to a linear equation.}  
  After taking logs, both equations become linear in $\log x$ and $\log y$, enabling elimination.

  \item \textbf{Factoring differences of squares.}  
  $(\log 2)^2 - (\log 3)^2 = (\log 2 - \log 3)(\log 2 + \log 3) \neq 0$, which justifies the division step.
\end{enumerate}
\end{propbox}

% =====================================================================
\section{Problem 2 — JEE Advanced 2018}
% =====================================================================

\begin{problembox}{Problem 2 \hfill \normalfont\small [JEE Advanced 2018]}
Find the value of
\[
E \;=\; \left[(\log_2 9)^2\right]^{\!\frac{1}{\log_2(\log_2 9)}}
\;\times\;
\left(\sqrt{7}\right)^{\!\frac{1}{\log_4 7}}
\]
\end{problembox}

\subsection*{Solution}

Define $E = A \times B$ where:
\[
A = \left[(\log_2 9)^2\right]^{\!\frac{1}{\log_2(\log_2 9)}}, \qquad
B = \left(\sqrt{7}\right)^{\!\frac{1}{\log_4 7}}
\]

\subsubsection*{Evaluating $A$}

Let $t = \log_2 9$. Then:
\[
A = (t^2)^{1/\log_2 t} = t^{2/\log_2 t}
\]

Now use the identity $t^{1/\log_2 t} = 2$ (since $\log_2 t \cdot \log_t 2 = 1$ implies $t^{1/\log_2 t} = 2$):
\[
A = t^{2/\log_2 t} = \left(t^{1/\log_2 t}\right)^2 = 2^2 = 4
\]

More explicitly, using $\log_2 9 = \log_2 9$:
\[
\frac{2 \log_2 9}{\log_2(\log_2 9)} \cdot \log_2(\log_2 9) \xrightarrow{\text{exponent}} (\log_2 9)^{2\log_2(\log_2 9)/\log_2(\log_2 9)} = (\log_2 9)^2
\]

\noindent Actually, let us proceed cleanly. Since $t^{1/\log_2 t} = 2^{\log_2 t / \log_2 t} = 2^1 = 2$:
\[
A = \left(t^2\right)^{1/\log_2 t} = t^{2/\log_2 t} = \left(t^{1/\log_2 t}\right)^2 = 2^2 = 4
\]

\subsubsection*{Evaluating $B$}

\[
B = (\sqrt{7})^{1/\log_4 7} = 7^{1/2 \cdot 1/\log_4 7}
\]

Convert $\log_4 7$ to base 2: $\log_4 7 = \dfrac{\log_2 7}{\log_2 4} = \dfrac{\log_2 7}{2}$.

So $\dfrac{1}{\log_4 7} = \dfrac{2}{\log_2 7}$. Thus:
\[
B = 7^{\frac{1}{2} \cdot \frac{2}{\log_2 7}} = 7^{1/\log_2 7}
\]

Now apply the same identity: $7^{1/\log_2 7} = 2^{\log_2 7/\log_2 7} = 2^1 = 2$.

Hence $B = 2$.

\subsubsection*{Combining}

\[
E = A \times B = 4 \times 2 = \boxed{8}
\]

\begin{answerbox}
\[E = 4 \times 2 = \mathbf{8}\]
\end{answerbox}

\begin{propbox}
\begin{enumerate}[leftmargin=*, label=\textbf{P\arabic*.}]
  \item \textbf{Key identity:} $a^{1/\log_b a} = b$.\\
  Proof: Let $\log_b a = k$, so $a = b^k$. Then $a^{1/k} = (b^k)^{1/k} = b$.

  \item \textbf{Change of base formula.}  
  $\log_4 7 = \dfrac{\log_2 7}{\log_2 4} = \dfrac{\log_2 7}{2}$.

  \item \textbf{Power rule of exponents.}  
  $(a^m)^n = a^{mn}$.

  \item \textbf{Decomposition.}  
  $\sqrt{7} = 7^{1/2}$, allowing the exponent to be combined before applying the identity.
\end{enumerate}
\end{propbox}

% =====================================================================
\section*{Summary of Key Log Properties}
% =====================================================================

\begin{mdframed}[backgroundcolor=boxbg, linecolor=problemblue, linewidth=1pt]
\textbf{Properties Consolidation}
\begin{align*}
&\bullet\ a^{1/\log_b a} = b && \text{(reciprocal log identity)}\\
&\bullet\ \log(ab) = \log a + \log b && \text{(product rule)}\\
&\bullet\ \log_b a = \frac{\log_c a}{\log_c b} && \text{(change of base)}\\
&\bullet\ (a^m)^n = a^{mn} && \text{(power of a power)}\\
&\bullet\ a^{\log_a b} = b && \text{(basic exponent-log inverse)}
\end{align*}
These five properties together are sufficient to solve the vast majority of JEE Advanced logarithm problems.
\end{mdframed}

\end{document}
